\documentclass[11pt,a4paper]{ltjsarticle}
\usepackage[margin=25mm]{geometry}
\usepackage{amsmath,amssymb,booktabs,longtable,array}
\usepackage[haranoaji]{luatexja-preset}
\usepackage[unicode,hidelinks]{hyperref}
\usepackage{xurl}
\title{nc-p1123研究方針と再開手順}
\author{}
\date{2026年11月19日\quad 第12版}
\begin{document}
\maketitle

\section{最初に読む：現在の目標と対象}
現在の主目標は、\textbf{整域性を仮定した、任意の$n\geq1$における
primitive jumps $(1^n,2,3)$の語長二次AS正則$\mathbb Z$-代数の構造}を明らかにすることである。
基礎体は標数零の代数閉体を基準とする。
正連結・局所有限、左右AS次元$n+2$、有限生成射影による単純加群の有限解消、
各行・各列の多項式増大を保つ。
最新の研究依頼では$\mathbb Z$-整域性を仮定してよい。
標準Hilbert級数、noetherian性、Koszul性、1周期性まで仮定してよいという意味にはしない。
quadraticは、重み1、2、3の最小生成元を全て語長1としたときの斉次二次性である。

この目標では、二段Ore表示を次の二つに分ける必要がある。
\begin{enumerate}
\item 単一のgraded代数$D=B[z;\sigma][w;\tau,\partial]$を用い、
$A\simeq\check D$とする表示。この表示は$A$の1周期性を含意する。
\item 生成元族$z_i,w_i$と添字に依存する交換写像・導分を用いる、
$\mathbb Z$-代数上の二段indexed Ore表示。これは1周期性を要求しない。
\end{enumerate}
\textbf{第一の表示が任意$n$の全対象で成立するという主張は、$n=1$の反例により偽である。}
第二の自然なindexed表示の網羅性は、$n=1,2$で証明した。
$n=1$では裸の添字保持同型による完全分類と1周期化の必要十分条件も得た。
\textbf{網羅性の未解決範囲は$n\geq3$である。}
$n=2$では単一のgraded二段Ore代数へ下降する、より強い構造定理が成立している。

正本は\texttt{uedakazushi/nc-p1123}の\texttt{main}。
再開時は本方針、\path{AGENTS.md}、\path{research/STATE.json}を最初に読む。
公表用原稿は\path{papers/nc-p1123.tex}、未収載の補助結果と障害は
\path{research/notes_ja.tex}、$n=2$の七窓表示は\path{research/indexed_z_algebra.tex}にある。
一次文献の適用範囲は\path{research/BASELINE.md}、実行した検証の範囲は
\path{research/validation.json}を参照する。

\section{$n=1$の完全分類と任意$n$への補助結果}
\subsection{$n=1$のindexed二段Ore表示と完全分類}
主論文\texttt{genidx:thm:one-light-ore}は、対象$\mathcal Q_1^{\mathrm{dom}}$の全てが
indexed二段Ore表示を持つことを証明する。
AS次元3と語長二次性からtopの語長3とKoszul性が従い、
整域性から$r_2=r_6=0$、AS対合からAS添字6と$r_3=r_4=r_5=1$が得られる。
topテンソルにより重み4関係の両混合係数が非零となり、
EulerのHilbert級数と整列語の生成性を比較してPBW基底と二段表示を得る。
Hilbert級数、noetherian性、Koszul性はここでの追加仮定ではない。

主論文\texttt{genidx:thm:one-light-classification}によれば、生成元を正規化すると
\begin{align*}
z_i x_{i+2}&=x_i z_{i+1},\\
w_i z_{i+3}&=c\,z_i w_{i+2},\\
w_i x_{i+3}&=x_iw_{i+1}+d_i z_i z_{i+2},
\end{align*}
\[
c\in k^\times,\qquad
d_{2j}=c^{-j}\alpha,\qquad d_{2j+1}=c^{-j}\beta
\quad(j\in\mathbb Z)
\]
と書ける。逆に、全ての$(c,\alpha,\beta)\in k^\times\times k^2$が
語長Koszul、左右noetherian、標準Hilbert級数を持つAS3 $\mathbb Z$-整域を与える。
この逆構成のAS性は具体的な左右解消によって証明した。

二つの正規形が裸の添字保持同型を持つための必要十分条件は、
$c'=c$かつ$(\alpha',\beta')=s(\alpha,\beta)$となる$s\in k^\times$が存在することである。
従って各$c$における同型類の\emph{集合}は
$\{0\}\sqcup\mathbb P^1(k)$である。これをaffine quotientやモジュライ空間の構成と同一視しない。
1周期となるための必要十分条件は
\[
\alpha^2=c\beta^2
\]
であり、零対も含む。それ以外の対象の最小正周期は2である。

\subsection{整域・noetherian・標準Hilbertでも1周期化しない例}
上の分類で$(c,\alpha,\beta)=(1,1,0)$を取る。
各始点$i$にjump 1、2、3の生成元$x_i,z_i,w_i$を置き、
\[
\epsilon_i=\begin{cases}1&i\text{ が偶数},\\0&i\text{ が奇数}\end{cases}
\]
とし、関係を
\begin{align*}
z_i x_{i+2}&=x_i z_{i+1},\\
w_i z_{i+3}&=z_i w_{i+2},\\
w_i x_{i+3}&=x_iw_{i+1}+\epsilon_i z_i z_{i+2}
\end{align*}
で定める。この$A$は、語長Koszul、AS次元3、AS添字6、左右noetherianな
$\mathbb Z$-整域であり、各行の二変数Hilbert級数は
\[
H_i(u,t)=\frac{1}{(1-ut)(1-ut^2)(1-ut^3)}
\]
である。最小正周期は2である。
完全な証明は主論文\texttt{genidx:thm:period-two-counterexample}にある。

語長二次PBW表示の唯一の重なり$wzx$は$\epsilon_i=\epsilon_{i+2}$によって解消する。
一方、重み4関係から$k(z_i z_{i+2})$への射影階数は偶数位置で1、奇数位置で0となる。
この階数列は、augmentation filtrationによる同型の線型化を用いると、
裸の添字保持同型でも保存される。そのため1周期同型は存在しない。
この反例は$n=1$の完全分類の一例であり、indexed表示の網羅性と両立する。

従って、標準Hilbert級数、Koszul性、noetherian性を追加しても、
単一のgraded二段Ore代数への全$n$での下降は回復しない。
$n=2$で得た2周期と7周期の結合を、任意$n$へ数字だけ置き換える方法も使えない。
例えば$n=1$ではAS添字6と2周期から1周期は従わない。

\subsection{任意$n$で使える確実な補助結果}
研究ノート\texttt{note:general-as-translation}では、有限AS解消と完全対合を持つ
語長二次対象について、AS頂点置換が$\ell(i)=i+L$という平行移動になり、
語長を保つ$A\simeq A(L)$が得られることを証明する。
$E^1$の矢印による到達可能性が整数の順序と一致するので、$n=1$も含む。
二次表示の復元には$E^1,E^2$とその積で足り、全ExtのKoszul生成性を仮定しない。
ただし、$L=n+5$や語長Koszul性まで得たとはしない。

研究ノート\texttt{note:general-normal-z}では、全始点で
\[
r_3(i)=n,\qquad r_5(i)=1
\]
が成立し、$z_i$が全て左右正則なら、$z_i$が正規元族をなし、
$A(2)\simeq A$となることを証明する。
これは$n=2$の正規元族の議論を任意$n$へ移したものである。
整域性からは関係射影の単射性と上界$r_3(i)\leq n$, $r_5(i)\leq1$, $r_6(i)=0$が出るが、
その上界の達成は$n=1,2$では証明したが、$n\geq3$では未証明である。

研究ノート\texttt{note:one-heavy-frobenius}は、重い生成元を一つ持つ
\emph{graded}代数に関するFrobenius対合の補助結果である。
$Q=T(V\oplus kw)/(R\oplus S)$、$R\subseteq V\otimes V$、
$S\subseteq Vw\oplus wV$とし、$Q$がKoszul AS次元$n+1$、
$B=T(V)/(R)$自身がKoszul AS次元$n$で、有限解消を持つことを仮定すると、
$Q\simeq B[w;\mu]$、$\mu\in\operatorname{Aut}_{\mathrm{gr}}(B)$を得る。
noetherian性、指定Hilbert級数、$Q$の整域性は仮定しない。
$B$自身のAS性などの仮定を省略したり、一般の$\mathbb Z$-代数に対する
二段表示の網羅定理と扱ったりしない。

\subsection{$n\geq3$への一般化に残る論理上の障害}
$n=1$では$E^1\times E^2\to E^3$からtopの語長3を得る。
$n=2$ではAS次元4の$E^2\times E^2\to E^4$の完全対合から、
topの語長が4となり、語長Koszul性を得た。
$n\geq3$では$E^2$の相補相手が$E^n$となるため、この議論はそのまま移らない。
既存の任意$n$のgraded証明は、指定Hilbert級数からAS添字$n+5$を導いた後に
差次数を用いてKoszul性を証明している。
従って、Hilbert級数を仮定しないindexed対象にその数値結論を引用しない。

また、基底となる純jump 1部分$B$のAS正則性と、
正規商での重み3生成元族の交換写像を別々に確立する必要がある。
元の$A$が整域であることから、商$A/(z)$が整域であると結論しない。
有限次のFrobenius性やHilbert一致だけで、全Koszul複体の正確性を認定しない。

\section{優先順位と次の具体的操作}
\paragraph{P1：$n\geq3$のindexed二段Ore表示の網羅性。}
$n=1$の完全分類と$n=2$の構造定理を出発点とし、$n\geq3$の整域対象について、
基底$\mathbb Z$-代数$B$、jump 2と3の生成元族、添字移動を伴う交換写像と導分を構成する。
単一のgraded代数への下降をindexed表示の定義に埋め込まない。
$r_3=n$, $r_5=1$の強制、基底$B$のAS性、正規商の第二段という
三つの段階に分け、仮定なしに証明できる範囲を調べる。
特定の中間仮定を置いた定理が得られた場合は、その仮定と全対象への適用の未証明部分を明記する。

\paragraph{P2：$n\geq3$での混合関係数と基底の正則性。}
整域性による$r_3\leq n$, $r_5\leq1$, $r_6=0$から出発する。
AS平行移動の存在は使えるが、その具体値と語長Koszul性を先取りしない。
最小解消の語長・重み・差次数を保って相補対合を追い、
混合関係数の不足が多項式増大と両立するか調べる。
標準Hilbert級数を仮定する補助問題と、仮定しない主問題を別に記録する。
基底$B$のAS性を確立できれば\texttt{note:one-heavy-frobenius}を利用する。

\paragraph{P3：周期構造とgraded代数への下降条件。}
$n=1$では$\alpha^2=c\beta^2$が完全な判定であり、$n=2$では全整域対象が下降する。
$n\geq3$で正しい条件を探す。
AS平行移動$L$と正規元族による2周期が得られれば、$\gcd(L,2)=1$のときは1周期となる。
ただし$L=n+5$は$n\geq3$の主問題で未証明である。
$L$が偶数の場合には、$\rho_i$などの添字保持不変量と基底変更方程式を使い、
2周期構造が1周期へ下降する条件を調べる。
既存の整域Oreデータの有限同型条件は、既にgraded表示を持つ対象の比較に使う。
それを全indexed対象の1周期化に使用しない。

\paragraph{P4：$n=2$の全分類を副課題として保持する。}
元の四次元対象$\mathcal Q$の全分類は未完成であり、成果は全て保持する。
現在は整域性が許されているため、その自動性を主課題の前提作業にはしない。
副課題として、七つの$z_i$の全次数の左右正則性、混合関係射影の退化、
最小周期7の零因子例の存在を調べる。
逆向き片側$R_5$関係の位置差2、3、4（法7）は
研究ノート\texttt{note:r5-orientation}で排除済みであり、残る差1、5、6と同方向だけの退化が対象となる。
整域Ore部分の一般軌道の標準形・安定化群・成分・特殊化も保持する。

語長二次性を外す$n\geq3$のgraded問題や、幾何学的認識はさらに副次的な課題とする。
次数1の窓の数だけを置き換えず、各jumpの始点・終点と語長を追う。

\section{保持する主要な結果と適用範囲}
\subsection{$n=2$の全対象と整域部分}
ここで$\mathcal Q$は、整域性を課さない重み$(1,1,2,3)$・AS次元4の従来の対象である。
全$\mathcal Q$について、語長Koszul性、AS添字7、7周期性、標準関係数とHilbert級数を証明した。
\[
\mathcal Q=\mathcal H=\mathcal S,\qquad
H_i(u,t)=\frac{1}{(1-ut)^2(1-ut^2)(1-ut^3)}.
\]
語長線型の最小解消の重みシフトは
\[
(0);\ (1,1,2,3);\ (2,3,3,4,4,5);\ (4,5,6,6);\ (7)
\]
である。これを任意$n$について証明したと読み替えない。

外側の関係数$r_2=r_5=1$は解析的に強制した。
混合関係数$r_3=r_4=2$は七周期Hilbert行列の成長から必要条件を導き、
890個の整数列の行列式を厳密に検査する計算補助証明である。
証明は主論文\texttt{idx:thm:outer-relation-ranks},
\texttt{idx:thm:mixed-relation-ranks}, \texttt{idx:cor:all-standard}にある。
コード\path{computations/verify_mixed_ranks.py}と
証明書\path{computations/mixed_ranks_certificate.json}を保持する。

$\mathcal Q$の整域部分では、1周期性と
\[
A\simeq\check{\bigl(B[z;\sigma][w;\tau,\partial]\bigr)}
\]
を証明した。$B$は量子平面またはJordan平面で、noetherian性も結論である。
証明は\texttt{idx:thm:domain-standard}から\texttt{idx:thm:domain-ore}にある。
全$\mathcal Q$について整域性は、選んだ表示の$z_0,\ldots,z_6$が
全次数で左右正則であることと同値である（\texttt{idx:prop:z-regular-criterion}）。
有限添字の判定であり、有限次数での検査とは区別する。

\subsection{同型判定と七窓}
既存の整域Oreデータについて、$S=\sigma|_V$, $M=\mu|_V$、
$T=MS^{-1}$、$J=c^{-1}M^2S^{-3}$、$\lambda=c/c'$、$P_i=T'^{-i}PT^i$と置く。
裸の添字保持同型の存在は、$P\in\mathrm{GL}_2(k)$、$\kappa\ne0$で
\[
J'P=PJ,\quad(P\otimes P_1)r\in k^\times r',\quad
L'P_3=\kappa L,\quad L'P_4=\kappa\lambda L
\]
を満たすことと同値である（\texttt{idx:thm:finite-isomorphism}）。
任意$n$の既存の二段Ore族にも、基底関係空間を対応させる形で適用できる。
AS正則性の有限検査でも、全非線型自己同型群の記述でもない。

正規化した対角三係数族の全同型は、一致または
\[
(a,b,c)\longmapsto(a^{-1},b^{-1},ca^3/b^2)
\]
による一致であり、全特殊値を含む。
固定点は$a=1$, $b=\pm1$である。
対角族の同型類の有限商の次元は3となる
（\texttt{diag:thm:indexed-isomorphism}, \texttt{diag:cor:quotient}）。
全ASモジュライの成分・次元と同一視しない。

$n=2$の全対象は固定階数の七窓表示に入り、表示空間の次元は98、
線型基底変更群の次元は42である。差56をモジュライ次元とはしない。
七窓内のAS/Koszul locus、軌道、零因子を持つ枝の幾何は未決である。
既知の対角族、非零導分を持つ初期Ore族、射影階数$\rho_i$は回帰例・不変量として保持する。

\subsection{任意$n$の既存graded定理}
複素数体上の連結・局所有限・左右noetherianなgraded AS正則整域について、
大域次元$n+2$、最小生成元次数$(1^n,2,3)$、語長二次性および
\[
H_A(t)=\frac{1}{(1-t)^n(1-t^2)(1-t^3)}
\]
を仮定すれば、二段Ore表示は網羅的であり、基底$B$はAS次元$n$の整域となる。
これは主論文\texttt{thm:general-quadratic-classification}の定理である。
この定理の仮定は、その定理を使う範囲では保持する。
最新のindexed主問題でnoetherian性やHilbert級数が許可されたという意味ではない。

同稿は一般の斉次Oreデータの有限係数条件、全graded同型、
語長二次部分の同型存在条件も含む。
$n=2$では語長二次性を仮定しないnoetherian・指定Hilbertのgraded分類もある。
最小生成元次数だけを固定すると大域次元が$n+m+2$となる明示族も保持する。
この次元無界の例は、大域次元$n+2$を固定した主問題の反例ではない。
またSklyanin基底を用いる例は、二段Oreの基底まで多項式型に分解できるという主張と区別する。

\section{達成基準・公表原稿・再現}
現在の任意$n$の課題では、反例が否定する範囲、正しいindexed表示の定義、
その必要十分条件または明示した仮定の下の網羅定理、graded代数への下降条件を区別して確立する。
$n=1$のindexed網羅性・完全分類・下降条件は証明済みである。
$n=2$のindexed網羅性とgraded下降も証明済みだが、その全対象の完全分類は未完成である。
一般indexed網羅性は$n\geq3$で未完成である。
反例や補助定理の証明を記載したこと、内部監査、計算、形式化、第三者査読、新規性確認は別の状態とする。
形式化・第三者査読・新規性の最終確認は未完である。

公表用原稿は\path{papers/nc-p1123.tex}に統合する。
研究ノートに置いた補助結果は、証明と適用仮定を保ち、完成した一般定理へ格上げしない。
文献PDF、検索結果の大量出力、キャッシュ、旧原稿、作業実況を一式へ収録しない。

\path{computations/weighted_baseline.py}は基準族のHilbert数値、重なり、対角twistを検算する。
\path{computations/verify_mixed_ranks.py}は$n=2$の890候補を再計算して証明書と照合する。
この計算を任意$n$の証明として使用しない。
\path{scripts/check_project.py}はファイル、参照、主張の所在を検査するもので、数学的証明や組版ではない。
再現コマンドとLuaLaTeX設定はREADME、実際の検証状態は\path{research/validation.json}に置く。

\section{常設の研究原則と保存許可}
本節は本研究固有の優先順位とは独立した継続規則である。
数学的重要性、新規性、理解の深化、応用可能性、正確性、実現可能性を総合し、
研究全体の価値が最大となるように目標と方法を柔軟に改める。
重要な変更は現在の理由と元の問題の解決・未解決を短く整理し、履歴として蓄積しない。

主張の仮定、依存、証明位置、未解決箇所を明示する。
証明済み、追加仮定付き、未完成、予想、計算観察を区別し、
誤りの訂正時は関係する本文・要約・方針を整合させる。
一次文献の内容・書誌を推測で補わず、未確認点と既知結果との差を記す。

論文は日本語LuaLaTeXで現在の成果を理解しやすい順序に編成する。
有効で公表価値のある証明を軽量化のために削除しない。
研究ノートは現在の知識、例、反例、障害、着想を残し、本文の証明を全文重複しない。
再開に必要な情報を会話だけに置かず、現行の方針・論文・ノート・必要計算で自己完結させる。

研究の区切りごとに、\textbf{ユーザーの承諾を改めて求めず現行の研究継続用一式を保存する}
常設許可がある。ソース正本はGitHub mainとし、約15分を目安に安全な区切りで直接pushする。
Git管理物はLibraryへ重複保存しない。独立した非追跡の配布成果は、
利用可能な保存先に常設許可の範囲で保存する。
Library上の既存ファイルの削除は別操作であり、明示指示なしには行わない。

一式には現行版だけを収め、旧ZIP・旧原稿・バックアップ・ログ・中間生成物を入れない。
日付と版を揃え、再生成手順と証明に必要な入力を保持する。
ファイルの生成、数学的検証、組版、遠隔保存を別状態として報告する。
失敗した組版や未実行の検証を完了扱いしない。
この運用は実行時間を延長せず、タスク終了後の自動再開を保証しない。
\end{document}
